% !TEX program = lualatex
% ============================================================
% main.tex — the collection itself
%
% This file is structure and poems. All typographic choices
% live in poetrystyle.sty, which also loads babel and sets
% the paper size.
%
% Compile with LuaLaTeX, not pdflatex:
%   latexmk -lualatex main
% ============================================================
\documentclass[12pt,twoside]{book}


% ============================================================
% LOAD POETRY STYLE
% ============================================================

% booklet is on in every variant: the collection is padded
% with blank pages to a page count divisible by four, so a
% folded booklet has a back. Remove the option for loose sheets.

% ------------------ Print, 140 x 216 mm -----------
\usepackage[booklet]{poetrystyle}

% ------------------ Printing at home, A5 ----------
% Same text block, so line and page breaks are as above.
% \usepackage[a5,booklet]{poetrystyle}

% ------------------ Centred poems -----------------
% Each poem centred on its longest line.
% \usepackage[centered,booklet]{poetrystyle}

% ------------------ Page proofs -------------------
% The page untouched, centred on A4 with crop marks.
% \usepackage[proof,booklet]{poetrystyle}


% ============================================================
% DOCUMENT
% ============================================================
\begin{document}

\title{Title of the Collection}
\subtitle{Poems}
\author{Author Name}
\date{\the\year}


% ============================================================
% FRONT MATTER
% ============================================================

% ------------------ Half-title --------------------
% Optional: the title alone on a page before the title page.
% \halftitlepage

% ------------------ Title page --------------------
\booktitlepage

% ------------------ Colophon ----------------------
\colophonpage{%
  © \the\year\ Author Name\\
  Set in Garamond Libre with \LaTeX\\
  ISBN 978-0-000-00000-0}

% ------------------ Dedication --------------------
\dedicationpage{To someone.}

% ------------------ Contents ----------------------
\poemtoc


% ============================================================
% THE POEMS
%
% One poem environment per poem. Every poem begins on a new
% page. Replace the placeholders below with your own.
% ============================================================

\part{I}

\begin{poem}[A poem with a title]
The title goes in the brackets
after begin.

Each line in the source
is a line on the page,
and a blank line
gives a new stanza.
\end{poem}

\begin{poem}
A poem without a title
begins where the title would have stood

and is listed in the contents
by its first line.
\end{poem}

\begin{poem}[Indentation]
Spaces at the start of a line
    indent the line,
        half an em for each,

and spaces   inside the line   are kept.
\end{poem}

\begin{poem}[A long line]
A line can be short.

A line that is longer than the text block is wide is broken by itself, and the continuation is indented so it cannot be mistaken for a new line.

The next line begins at the left edge again.
\end{poem}


\part{II}

\begin{poem}[A long poem]
This poem is longer than the page.
It is therefore divided in two,
but never in the middle of a stanza:
the break falls in the space between two of them.

The stanza that does not fit
is moved whole to the next page,
and then the bottom of the page
may be left open.

This is what the third stanza looks like.
It has four lines like the others,
and does not say much,
but it takes the room it should.

The fourth stanza follows.
The lines stand sixteen points apart,
and between the stanzas
there is one line of space.

The fifth stanza is the last on this page.
The text block holds twenty-nine lines,
and the poem begins on the sixth,
with the title on the fourth.

The sixth stanza stands at the top of the next page.
The continuation begins at the top of the text block,
without the space that stands above the title,
since it is not a new poem.

The page number is centred at the bottom,
one step smaller than the text,
and there is no running head
above the poems.

The eighth stanza is the last.
The next poem begins on a new page,
at the same height as this one began,
and so the collection goes on.
\end{poem}

\begin{poem}
Replace these poems with your own.

That is all it takes.
\end{poem}


\end{document}
